\section{Evaluation}\label{sec:experiment}

\subsection{Two Legal Retrieval Regimes}

BarExamQA and HousingQA sit at opposite ends of the question-corpus
vocabulary gap. BarExamQA questions are fact patterns in narrative prose,
while the support corpus uses doctrinal vocabulary --- a large gap that a
generated legal passage can close. HousingQA questions are
statutory-entailment questions framed in landlord-tenant terms, sharing
vocabulary with the housing statutes that answer them --- a small gap where
raw queries are already corpus-shaped. The main paper reports both
benchmarks, exact-scored.

\paragraph{BarExamQA.}
BarExamQA contains 1,195 multiple-choice bar-exam questions in the evaluation
split used here. The retrieval corpus is a frozen support-passage collection
with official gold passage ids. Hit@5 and MRR@5 therefore measure whether the
retriever exposes the labeled rule passage before the answer model chooses an
option.

\paragraph{HousingQA.}
HousingQA contains 6,853 yes/no statutory entailment questions over state
housing statutes. The corpus contains many jurisdictions, and each question
names the relevant state. All main HousingQA retrieval rows apply that state
as a metadata filter before ranking. Rows that do not use this filter are
excluded from the main HousingQA comparison.

\subsection{Models, Baselines, and Retrieval Stack}

We evaluate three model configurations: Llama 3.1 8B Instant, Gemma 4 26B,
and Llama 3.3 70B Versatile. Within each row the same model serves as both
the query generator (when one is needed) and the final answer model, so
accuracy lift cannot be attributed to a stronger reader.

We compare four deployable retrieval methods and one oracle control. All
non-oracle rows answer from $k=5$ retrieved passages.

\begin{itemize}
\item \textbf{LLM only} answers the original question without retrieval.
\item \textbf{Raw question RAG} retrieves with the question text, then
answers from the question plus top-$k$ evidence.
\item \textbf{HyDE} retrieves with a generated legal passage, without a
draft answer prior.
\item \textbf{\hyre{} (ours)} retrieves with a draft-guided legal search
passage; the draft answer is private and is not shown to the final answer
call.
\item \textbf{Gold passage} supplies the labeled passage directly as an
oracle control for evidence use (not a deployable retrieval method).
\end{itemize}

Structured rewrite and gold-plus-neighbor diagnostics appear in the
appendix; route labels are retained in the released artifacts for
reproducibility.

\paragraph{Retrieval stack.}
Runs use seed 42, temperature 0, and a 2048-token answer cap. If the model
produces a malformed generated passage or no parseable final answer, the row
fails rather than silently falling back to raw-question retrieval. Retrieval
uses Chroma collections built with \method{Alibaba-NLP/gte-large-en-v1.5}
embeddings; a cross-encoder (\method{ms-marco-MiniLM-L-6-v2}) then reranks
the candidate set before the top-$k$ cut.

\paragraph{Generated-query evaluation.}
Generated-query rows are evaluated in two stages. First, the query-generation
call writes the HyDE passage or \hyre{} search passage. Then retrieval stores
the ranked document ids. The answer run replays those ids strictly; an
invalid generated passage or missing cache entry is not replaced with
raw-question retrieval. This keeps the comparison about the method's actual
search text.

\paragraph{Per-dataset retrieval interfaces.}
BarExamQA retrieval uses the shared fact pattern and question stem;
generated methods may see option text but no correctness marker. HousingQA
retrieval applies the state metadata filter described above. The filter is
not a \hyre{}-specific feature: it is the benchmark interface used by every
HousingQA row in the main paper.

\subsection{Answer and Evidence Metrics}

The downstream metric is exact answer accuracy. BarExamQA uses the last valid
multiple-choice answer line, and HousingQA uses \texttt{Answer: Yes} or
\texttt{Answer: No}. Retrieval is measured by Hit@5 and MRR@5: whether any
labeled support passage appears in the top five, and how early it appears.

We also report efficiency as the number of model calls per question, the
tokens the answer call processes per question (the ``answer-stage tokens''),
and correct answers per million answer-stage tokens. Two-call methods are
not penalized in this count: the first-stage query-generation call is
counted as part of the method footprint and logged separately. The
answer-stage token count shows how much context the final answer call sees
once retrieval has selected evidence.

The main tables include only validated full rows. Missing cells are not
imputed. Appendix~\ref{sec:appendix:policy} lists the inclusion rule,
coverage grid, supplemental controls, top-$k$ diagnostics, efficiency table,
and reproducibility details.
